% ==========================================================================
% CHAPTER 5 -- ADDRESSING PROGRAM OUTCOMES, COMPLEX ENGINEERING PROBLEMS
%              AND ACTIVITIES
%
% The three tables contain the OFFICIAL statements copied from Annex V of
% the EEE 4000 manual. The right-hand column of each table points to the
% specific sections of this book where each outcome/attribute is addressed.
% ==========================================================================

\chapter{Addressing Program Outcomes, Complex Engineering Problems and Activities}
\label{ch:po}

This chapter maps the work presented in Chapters~\ref{ch:introduction}--\ref{ch:social}
onto the program outcomes (POs), complex engineering problem (CEP)
attributes, and complex engineering activity (CEA) attributes defined in
the EEE 4000 manual. For each item, the section of this book that
demonstrates the outcome is identified, together with a one-line
explanation of how the work addresses it. The statements in the tables are
quoted verbatim from Annex V of the manual.


\section{Addressing Program Outcomes}
\label{sec:po}

Table~\ref{tab:po_mapping} maps the twelve program outcomes of the EEE 4000
manual onto the specific evidence in this book. The work is a
computational, research-based thesis, so the strongest outcomes are
PO(a)--PO(g) and PO(l); teamwork and management outcomes (PO(i), PO(j),
PO(k)) are addressed through the project execution process itself.
\begin{longtable}{@{}>{\raggedright\arraybackslash}p{0.55in}>{\raggedright\arraybackslash}p{3.1in}>{\raggedright\arraybackslash}p{1.5in}@{}}
\caption{Program Outcomes (POs) of the Project \& Thesis (EEE 4000)}
\label{tab:po_mapping} \\
\toprule
\textbf{PO} & \textbf{Statement} & \textbf{Addressing the PO} \\
\midrule
\endfirsthead
\multicolumn{3}{@{}l}{\small\textit{Table \thetable\ continued from the previous page}} \\
\toprule
\textbf{PO} & \textbf{Statement} & \textbf{Addressing the PO} \\
\midrule
\endhead
\bottomrule
\endfoot

PO(a) & Apply knowledge of mathematics, natural science, engineering
fundamentals and an engineering specialization as specified in K1 to K4
respectively to the solution of complex engineering problems.
(Engineering knowledge) & Sections~\ref{sec:data_analysis}
and~\ref{sec:math}: signal processing (Butterworth
filtering, Pan--Tompkins), neural network mathematics, and the statistical
formulation of the classifiers. \\
\addlinespace

PO(b) & Identify, formulate, research literature and analyze complex
engineering problems reaching substantiated conclusions using first
principles of mathematics, natural sciences and engineering sciences.
(K1 to K4) (Problem analysis) &
Sections~\ref{sec:literature}--\ref{sec:gap}: literature review and
research gap identify the protocol and minority-class problems; the
ablation study (Section~\ref{subsec:ablation}) reaches substantiated
conclusions. \\
\addlinespace

PO(c) & Design solutions for complex engineering problems and design
systems, components or processes that meet specified needs with appropriate
consideration for public health and safety, cultural, societal, and
environmental considerations. (K5) (Design/Development of solutions)
& Sections~\ref{sec:math}--\ref{sec:proposed_methodology}: design of the
hybrid and MSFT-Net architectures; Chapter~\ref{ch:social} addresses the
health, safety, legal, and societal considerations of the design. \\
\addlinespace

PO(d) & Conduct investigations of complex problems using research-based
knowledge (K8) and research methods including design of experiments,
analysis and interpretation of data, and synthesis of information to provide
valid conclusions. (Investigation)
& Chapter 3: controlled ablation experiments, paired McNemar tests,
Wilson and bootstrap confidence intervals, and the interpretation of the
results. \\
\addlinespace

PO(e) & Create, select and apply appropriate techniques, resources, and
modern engineering and IT tools, including prediction and modelling, to
complex engineering problems, with an understanding of the limitations. (K6)
(Modern tool usage)
& Sections~\ref{sec:materials},
\ref{sec:proposed_methodology}--\ref{sec:workflow}: Python,
TensorFlow/Keras, scikit-learn, TFLite
Micro, and STM32 toolchains; the limitations are stated explicitly in
Sections~\ref{subsec:explain} (explainability) and~\ref{sec:limitations}. \\
\addlinespace

PO(f) & Apply reasoning informed by contextual knowledge to assess societal,
health, safety, legal and cultural issues and the consequent
responsibilities relevant to professional engineering practice and solutions
to complex engineering problems. (K7) (The engineer and the society)
& Chapter~\ref{ch:social}: financial and health influences, safety,
legal/standards and
cultural/ethical assessment of the wearable classifier. \\
\addlinespace

PO(g) & Understand and evaluate the sustainability and impact of
professional engineering work in the solution of complex engineering
problems in societal and environmental contexts. (K7)
(Environment and sustainability)
& Sections~\ref{sec:environment} and~\ref{sec:sustainability}:
environmental life-cycle assessment and
sustainability analysis of the device, linked to the SDGs. \\
\addlinespace

PO(h) & Apply ethical principles and commit to professional ethics and
responsibilities and norms of engineering practice. (K7) (Ethics)
& Section~\ref{subsec:safety_legal_cultural}: honest reporting of the
near-zero faithfulness of the
attribution maps; Section~\ref{subsec:explain} and the negative results in
Sections~\ref{subsec:intra_compress} and~\ref{subsec:fbeat} are reported
rather than suppressed. \\
\addlinespace

PO(i) & Function effectively as an individual, and as a member or leader in
diverse teams and in multi-disciplinary settings.
(Individual work and teamwork) & The thesis was executed as a two-member
team: the dataset preparation and evaluation pipeline (Chapter~\ref{ch:methodology}) and the
deployment work (Section~\ref{subsec:mcu}) were divided and integrated; weekly
supervisor meetings required individual reporting and joint review. \\
\addlinespace

PO(j) & Communicate effectively on complex engineering activities with the
engineering community and with society at large, such as being able to
comprehend and write effective reports and design documentation, make
effective presentations, and give and receive clear instructions.
(Communication) & This book itself, the accompanying defence presentation,
and the manuscript prepared for journal submission communicate the work to
the engineering community. \\
\addlinespace

PO(k) & Demonstrate knowledge and understanding of engineering management
principles and economic decision-making and apply these to one's own work,
as a member and leader in a team, to manage projects and in
multidisciplinary environments. (Project management and finance)
& The project was scheduled across the EEE 4000 timeline with milestone
reviews; the budget was constrained to free public data, open-source tools,
and low-cost commodity hardware, as justified in
Section~\ref{subsec:financial_health}. \\
\addlinespace

PO(l) & Recognize the need for, and have the preparation and ability to
engage in independent and life-long learning in the broadest context of
technological change. (Lifelong learning)
& Sections~\ref{sec:literature} and~\ref{sec:future}: the literature review
and the future plan show
engagement with the rapidly evolving fields of deep learning and TinyML,
and identify the specific next steps to continue learning. \\

\end{longtable}

The mapping above is substantiated in prose as follows. The engineering
knowledge outcome (PO(a)) is demonstrated by the full pipeline: the
Butterworth filtering and Pan--Tompkins R-peak detection of
Section~\ref{sec:data_analysis} apply classical signal-processing theory,
while Equations~\ref{eq:attention}--\ref{eq:film} of
Section~\ref{sec:math} apply the mathematics of attention and gating to a
biomedical signal. Problem analysis (PO(b)) is demonstrated throughout
Chapter~\ref{ch:results}: the study is designed as a controlled experiment
with explicit partitions, the hypotheses are tested with McNemar's test and
Holm correction (Section~\ref{sec:eval_criteria}), and the negative results
(the collapsed $A2$ augmentation arm in Table~\ref{tab:fablation}, the
failed dynamic-range quantisation in Section~\ref{subsec:intra_compress})
are reported and explained rather than hidden. Design and development
(PO(c)) is demonstrated by the MSFT-Net architecture itself, whose design
decisions (multi-scale depthwise-separable convolutions, SE gating, FiLM
conditioning) are each justified in Sections~\ref{sec:math}
and~\ref{subsec:ablation} by ablation evidence rather than assertion.

Investigation (PO(d)) has the strongest supporting evidence: the twelve-arm ablation
study of Section~\ref{subsec:ablation}, the fusion-beat staged
intervention of Section~\ref{subsec:fbeat}, and the cross-protocol
synthesis of Section~\ref{sec:crossprotocol} are genuine investigations
that produce statistically supported conclusions. Modern tool usage (PO(e))
is shown by the use of state-of-the-art deep learning frameworks and the
embedded deployment toolchain, with the limitations of each acknowledged:
the explainability maps fail the faithfulness test (Section~\ref{subsec:explain}),
and the host-measured latency of the tiny model is distinguished from the
on-device measurement (Sections~\ref{subsec:mcu}). The engineer-and-society
(PO(f)) and environment-and-sustainability (PO(g)) outcomes are addressed
in Chapter~\ref{ch:social}, which assesses health, financial, legal, and
environmental impacts in the Bangladeshi context. Ethics (PO(h)) is
demonstrated by the candid reporting policy of this book: negative results
are reported as such, the $Q$-class caveats are stated explicitly
(Section~\ref{subsec:overall}), and no claim of clinical validity is made
beyond what the data support.

The teamwork (PO(i)), communication (PO(j)), and management (PO(k))
outcomes are addressed through the execution of the project itself rather
than through any single technical section. Finally, lifelong learning
(PO(l)) is demonstrated by the literature review of Section~\ref{sec:literature},
which spans 2015--2026 work across six research directions, and by the
concrete future plan of Section~\ref{sec:future}.


\section{Addressing Complex Engineering Problems}
\label{sec:cep}

The thesis addresses a complex engineering problem because it involves
research-based knowledge at the forefront of the discipline, wide-ranging
technical issues, and non-obvious design decisions. Table~\ref{tab:cep_mapping}
justifies each attribute of a complex engineering problem with evidence from
this work.

\begin{longtable}{@{}>{\raggedright\arraybackslash}p{1.4in}>{\raggedright\arraybackslash}p{2.4in}>{\raggedright\arraybackslash}p{1.4in}@{}}
\caption{Complex Engineering Problem Solving through the Project \& Thesis (EEE 4000)}
\label{tab:cep_mapping} \\
\toprule
\textbf{Attribute} & \textbf{Statement} & \textbf{How It Is Addressed} \\
\midrule
\endfirsthead
\multicolumn{3}{@{}l}{\small\textit{Table \thetable\ continued from the previous page}} \\
\toprule
\textbf{Attribute} & \textbf{Statement} & \textbf{How It Is Addressed} \\
\midrule
\endhead
\bottomrule
\endfoot

Depth of knowledge required (P1) & Requires research-based knowledge, much
of which is at, or informed by, the forefront of the professional discipline
and that allows a fundamental-based, first principles analytical approach
& Section~\ref{sec:literature}: the solution draws on 2015--2026
state-of-the-art research
(attention, SE, FiLM, TinyML); Section~\ref{sec:math} derives the
architectures from first principles. \\
\addlinespace

Range of conflicting requirements (P2) & Involve wide-ranging or conflicting
technical, engineering and other issues
& The accuracy--footprint--latency trade-off of Section~\ref{subsec:mcu},
and the conflict between permissive and record-disjoint evaluation
(Section~\ref{sec:crossprotocol}), are conflicting requirements resolved by
measurement. \\
\addlinespace

Depth of analysis required (P3) & Have no obvious solution and require
abstract thinking, originality in analysis to formulate suitable models
& Section~\ref{sec:proposed_methodology}: the fusion-beat treatment and the
auxiliary feature vector
are original formulations; Section~\ref{subsec:fbeat} shows the naive
solution ($A2$)
fails and only a staged, analysed design works. \\
\addlinespace

Familiarity of issues (P4) & Involve infrequently encountered issues
& The $S$ and $F$ classes (Sections~\ref{sec:background}
and~\ref{subsec:fbeat}) are rare beats whose
correct classification is unfamiliar to generic models; the $Q$-class
single-patient caveat (Section~\ref{subsec:overall}) is an infrequently
encountered statistical issue. \\
\addlinespace

Extent of applicable codes (P5) & Are outside problems encompassed by
standards and codes of practice for professional engineering
& The ANSI/AAMI EC57 standard (Section~\ref{sec:eval_criteria}) governs
reporting but does not
prescribe an architecture or an imbalance strategy;
Sections~\ref{subsec:ec57}--\ref{subsec:ablation}
show the standard must be combined with novel design work. \\
\addlinespace

Extent of stakeholder involvement and conflicting requirements (P6) &
Involve diverse groups of stakeholders with widely varying needs
& Section~\ref{sec:societal}: patients (early detection), cardiologists
(alarm burden,
trust in explanations), and low-income users (cost) have different needs;
Section~\ref{subsec:roc} quantifies the trade-off between detection and
false alarms. \\
\addlinespace

Interdependence (P7) & Are high-level problems including many component
parts or sub-problems
& The pipeline integrates preprocessing, feature engineering, two
architectures, imbalance handling, explainability, and deployment; the
ablation of Section~\ref{subsec:ablation} shows the components are
interdependent, with
each block critical for a different class. \\

\end{longtable}


\section{Addressing Complex Engineering Activities}
\label{sec:cea}

The development of this system also constitutes a complex engineering
activity. Table~\ref{tab:cea_mapping} justifies each attribute, drawing on
Chapter~\ref{ch:social} for the societal and environmental consequences
(A4).

\begin{longtable}{@{}>{\raggedright\arraybackslash}p{1.4in}>{\raggedright\arraybackslash}p{2.4in}>{\raggedright\arraybackslash}p{1.4in}@{}}
\caption{Complex Engineering Activities Solving through the Project \& Thesis (EEE 4000)}
\label{tab:cea_mapping} \\
\toprule
\textbf{Attribute} & \textbf{Statement} & \textbf{How It Is Addressed} \\
\midrule
\endfirsthead
\multicolumn{3}{@{}l}{\small\textit{Table \thetable\ continued from the previous page}} \\
\toprule
\textbf{Attribute} & \textbf{Statement} & \textbf{How It Is Addressed} \\
\midrule
\endhead
\bottomrule
\endfoot

Range of resources (A1) & Involve the use of diverse resources (and for this
purpose resources include people, money, equipment, materials, information
and technologies) &
Sections~\ref{sec:materials}--\ref{sec:workflow}: three public databases, GPU
training hardware, STM32 target hardware, Python/TensorFlow/TFLite
software, and the two authors' and supervisor's expertise. \\
\addlinespace

Level of interaction (A2) & Require resolution of significant problems
arising from interactions between wide-ranging or conflicting technical,
engineering or other issues
& Sections~\ref{subsec:intra_compress} and~\ref{subsec:mcu}: quantisation
interacts with accuracy, the
recurrent stage interacts with the evaluation protocol, and augmentation
interacts with the fusion class; each interaction was resolved by
experiment. \\
\addlinespace

Innovation (A3) & Involve creative use of engineering principles and
research based knowledge in novel ways
& Section~\ref{sec:proposed_methodology}: the staged fusion-beat pipeline
and the
MSFT-Net design (multi-scale depthwise-separable stem with SE and FiLM)
are novel combinations of published principles, validated in
Sections~\ref{subsec:fbeat}--\ref{subsec:ablation}. \\
\addlinespace

Consequences for society and the environment (A4) & Have significant
consequences in a range of contexts, characterized by difficulty of
prediction and mitigation
& Chapter~\ref{ch:social}: the consequences for cardiac patients and the
Bangladeshi
healthcare system are significant; the risk of misdiagnosis and alarm
fatigue is difficult to predict and is mitigated by statistical operating
point analysis (Section~\ref{subsec:roc}). \\
\addlinespace

Familiarity (A5) & Can extend beyond previous experiences by applying
principles-based approaches
& The work extends standard CNN/LSTM practice to the record-disjoint
protocol and to real microcontroller deployment, both of which go beyond
the experiences reported in most of the surveyed literature
(Section~\ref{sec:literature}). \\

\end{longtable}